\documentclass[12pt]{article}

\usepackage{url,verbatim, amsmath, amssymb, amsfonts, dsfont, color}
\textwidth 15true cm
\textheight 9true in
\oddsidemargin 0.30true in
\evensidemargin 0.30true in
\topmargin -0.3true in
\headsep 0.4true in

\def\sgn{{\rm sign}}
\def\drightarrow{{\stackrel{d}{\rightarrow}}}
\def\dotsim{\stackrel{.}{\sim}}
\def\goesd{\stackrel{d}{\rightarrow}}
\def\goesp{\stackrel{p}{\rightarrow}}
\def\goeslaw{\stackrel{{\cal L}}{\longrightarrow}}
%
\def\vp{{\varphi}}
\def\E{{\rm E}}
\def\Pvalue{{$p$-value}}
\def\Pvalues{{$p$-values}}
\def\pr{{\rm pr}}
\def\Pr{{\rm pr}}
\def\logit{{\rm logit}}
\def\T{{ \mathrm{\scriptscriptstyle T} }}
\def\diag{{\rm diag}}  
\def\half{{\textstyle{1\over 2}}}
\def\ith{{$i{\rm th}$ }}
\def\var{{\rm var}}
\def\cl{{{\rm c}\ell}}
\def\checkx{{\checked\hspace{-6pt}\setminus}}
\def\grad{{\nabla}}
\def\bs#1{{\boldsymbol{#1}}}

\sloppy %permits line-breaking of urls

\def\blue{\color{blue}}
\def\red{\color{red}}
\def\sgn{{\rm sign}}
\def\drightarrow{{\stackrel{d}{\rightarrow}}}
\begin{document}

\noindent{\bf Questions Week 6} {\hfill STA 2212S 2026}\\

%{\bf Due January 14} \hfill{\it MS, Exercise 5.2}

\bigskip

\begin{enumerate}

\item {\it SM Exercise 7.3.8} If $I$ is Bernoulli with $\Pr(I=1) = (\alpha_2-\alpha)/(\alpha_2-\alpha_1)$, and ${\cal R}_{\alpha_1}, {\cal R}_{\alpha_2}$ are rejection regions of size $\alpha_1, \alpha_2$, show that the rejection region $${\cal R} = I{\cal R}_{\alpha_1} + (1-I){\cal R}_{\alpha_2}$$ has size $\alpha$. {\it Note: this is the basis for randomized tests for discrete distributions.}
\item {\em  A Bayesian approach to  multiple testing} (adapted from Efron and Hastie \S 15.3)\\
Suppose we have computed $m$ test statistics $t_1, \dots, t_m$, and for each $j$ we will reject $H_{0j}$ if $t_j > t_0$, where $t_0$ is some cut-off to be determined. A statistical model for this could start with prior probabilities $\pi_0$ and $\pi_1 = 1-\pi_0$ that any $H_{0j}$ is true or not; and densities for $t_j$ either $f_0(t)$ or $f_1(t)$,  when $H_{0j}$ is true, or false, respectively. Since we don't know which $H_{0j}$ are true, the marginal density of $t_j$ is
\begin{equation}
f(t) = \pi_0 f_0(t) + \pi_1f_1(t).
\end{equation}
\begin{description}
\item{(i)} The {\em Bayes false-discovery rate} for cutoff $t_0$ is defined as
$$
Fdr(t_0) = \pr\{ H_{0j} \text{ is true } \mid t_j \ge t_0\}.
$$
Show that this is equal to $$\pi_0\{1-F_0(t_0)\}/\{1-F(t_0)\},$$ 
where $F(t)$ is the cumulative distribution function for the density (1). 
\item{(iii)} We don't know $F(t_0)$, but a reasonable estimate is the empirical c.d.f   $\widehat F(t_0) = m^{-1}\sum_{j=1}^m 1\{t_j \le t_0\}$ and hence
$$
\widehat{Fdr}(t_0)=\pi_0\{1-F_0(t_0)\}/\{1-\widehat F(t_0)\},
$$  Show that for $0 < q < 1$, the condition $$p_{(i)}\le (i/m)q$$ is equivalent to $$\widehat{Fdr}(t_{(i)}) \le \pi_0q.$$ 
\end{description}

\item State and prove the Neyman-Pearson Lemma. 

\textcolor{red}{
\textbf{Lemma:} If we consider parameter space $\Theta = \{\theta_0, \theta_1\}$ and a test $H_0: \theta = \theta_0$, $H_1 : \theta = \theta_1$, then among all tests of size at most $\alpha$, the likelihood ratio test, defined with test statistic 
\begin{equation*}
    T = \frac{f(x ; \theta_1)}{f(x ; \theta_0)} = \frac{f_1(x)}{f_0(x)}
\end{equation*}
with rejection region $R = \{T(X) \geq \kappa_\alpha\}$ where $\alpha \geq P_{H_0}(T(X) \geq \kappa_\alpha)$ is the most powerful test. \\
\textit{Proof.} Let $\tilde R$ be another rejection region with size at most $\alpha$. 
}

\end{enumerate}
\end{document}

